\documentclass{article}
\usepackage{url}
\usepackage{graphicx} % Required for inserting images

\usepackage{biblatex}
\addbibresource{./references.bib}

\title{Anycast and Third-party libraries IEEE Communications}
\author{
Hugo Pascual Adán
\and
Jose M. del Álamo
\and
David Rodriguez
\and
Juan C. Dueñas
}
\date{Sep 2025}

\begin{document}

\maketitle

\section{Abstract}

Third-party libraries offer a wide range of services that simplify, enhance, and accelerate the development of mobile apps. These libraries oftentimes rely on distributed infrastructures and technologies, such as anycast addresses, to deliver their services worldwide. Yet, they are also frequently involved in personal data transfers, which may cross international borders due to anycast working. However, major privacy legislations restrict the international transfers of personal data and set strict requirements to allow them. This article presents a novel approach to analyze the privacy compliance risks induced by third-party libraries using anycast communications. To this end, we conducted an experiment where we analyzed the personal data transfers of more than 5,759 Android apps, further identifying the libraries triggering the transfers and the destinations' geolocation. The results show that 90 percent of third-party libraries and 98.65 percent of apps integrating them potentially fail to meet the European Union requirements for international personal data transfers.

\section{Introduction}

Nowadays, the professional, cost-effective development of mobile apps heavily relies on third-party libraries (TPLs). TPLs offer pre-written code and functionalities that developers can easily integrate into their apps. Common uses for TPLs include social media login (e.g., logging in with Facebook or Google), analytics tools to track user behavior and app performance, or secure in-app payments and transactions. A significant amount of TPLs offer their services in a free or subsidized model, allowing the TPL provider to collect and utilize user data in return \cite{Zhan2022ResearchReview}.

The most integrated TPLs in mobile apps are provided by companies that operate globally, offering their services worldwide. Often, they leverage Internet Protocol (IP) anycast addresses \cite{Abley2006OperationServices} to ensure a similar user experience regardless of their users' location. In anycast, the same IP address is assigned to different network devices (typically servers) so that the network routes the IP packets to the one “closer” to the sender, in network terms. This allows service providers to improve service performance through lower latency and better resilience against attacks. Anycast addressing is widely used on the Internet, for example, in DNS root servers or cloud computing services such as Amazon Web Services and Microsoft Azure.

However, anycast great virtue can become a serious regulatory concern if personal data are transferred. Privacy and data protection regulations around the world have set strict requirements for international transfers of personal data. Regions and countries such as the European Economic Area, the United Kingdom, China, Australia, and the United States, among many others, regulate this aspect of data protection specifically (see the second section, where we provide an overview of the requirements set by different privacy and data protection legislation). Yet, anycast communications do not respect international borders as they send the data to the anycast instance closest to the source of the communication, in networking terms. As a result, apps integrating TPLs without being aware of their networking behavior could compromise their legal compliance with data protection regulations by transferring personal data to other countries. Furthermore, data subjects may lose the data protection given by their home legislation when their personal data crosses international borders, thus increasing the risk of leakages to unknown actors due to bad practices or insufficient protection and being subject to potential targeted privacy attacks like identity theft.

This article addresses the privacy compliance issues introduced by TPLs relying on anycast communications. To highlight the extent of the problem with a specific but relevant example, we conducted an experiment with 5,500 Android apps. We observed their behavior, traced their personal data transfers to anycast addresses, retrieved their privacy policies, and compared this information with the requirements set by the European Union's data protection legislation. This novel analysis revealed an alarming level of potential non-compliance issues. We further identified the TPLs initiating the transfers in each app and analyzed their privacy policies to understand whether they disclosed the international transfers. Considering these findings, we discuss the responsibility of the main agents involved in the Android development ecosystem, namely the app developers, TPL providers, and the marketplace owners. To our knowledge, this is the first study revealing the privacy risks introduced by TPLs relying on anycast communications.

\begin{figure}
    \centering
    \includegraphics[width=\textwidth]{statics/analysis_workflow.png}
    \caption{Workflow for the compliance analysis of personal data transfers to anycast addresses in Android apps.}
    \label{fig:analysis_workflow}
\end{figure}

\section{Legal Requirements for Personal Data Transfers}

Different legislations worldwide set requirements to allow transfers of personal data to third countries, often called international transfers. Some reasons for restricting and conditioning the sending of personal data abroad are to ensure that citizens' rights will be respected and to protect against external threats. Among others, the European Economic Area (EEA) with the General Data Protection Regulation (GDPR) \cite{REGULATIONSRelevance}, China with its Personal Information Protection Law (PIPL) \cite{TheTranslation}, Australia with the Australian Privacy Principles (APP) \cite{ReadOAIC}, the United Kingdom with the GDPR and the Data Protection Act \cite{Data2018} or the United States with the recent Protecting Americans' Data from Foreign Adversaries Act \cite{Rep.Pallone2024Text2024} are some important laws that restrict international transfers setting out the criteria and mechanisms through which they can be lawfully conducted.

Transparency and adequacy are two common requirements in most privacy legislation. Transparency requires the entities who transfer personal data to provide data subjects with clear, specific, and comprehensible information about the transfer. The guidelines issued by data protection authorities clarifying the information to be disclosed emphasize the need for data handlers to identify the destination country or region. For example, according to the Australian APP Article 5.2, the data handlers must notify individuals that their information may be disclosed to overseas recipients and the countries where they are likely to be located. The European, Chinese, and British legislations mandate similar requirements.

Adequacy requirements seek to ensure that the international recipient of personal data meets the domestic requirements or provides equivalent protection. Some countries (e.g., the United Kingdom or the European Union) have compiled a list of third countries deemed adequate, thus allowing the free transfers of personal data to them and mandating additional guarantees to those lacking this recognition. For example, the European Commission has issued adequacy decisions for 14 countries, including the United States, which has entered and exited the list on several occasions. On the other hand, the United States has recently enacted the Protecting Americans' Data from Foreign Adversaries Act \cite{Rep.Pallone2024Text2024}, which explicitly prohibits the transfer of its citizens' personal data to four specific countries: China, Iran, North Korea, and Russia. These new measures change previous US policies and, citing national security concerns, begin to regulate the sending of personal data to third countries or organizations through the implementation of a blacklist.

We have chosen the European Union GDPR as a reference to assess the compliance issues introduced by TPLs relying on anycast communications, as it has significantly inspired and influenced other privacy regulations worldwide \cite{Linden2020TheGDPR}. The GDPR imposes strict transparency requirements for the information to be provided to data subjects when their personal data is to be transferred to third countries, including the intention to carry out international transfers, the existence of an adequacy decision for the destination country and, otherwise, the additional safeguards put in place to ensure data protection. Furthermore, the guidelines of transparency issued by the European Data Protection Board explicitly state the need to name the third country to which data will be transferred \cite{JUSTICEwp260rev.01}.

\section{Analyzing Anycast Personal Data Transfers}

Figure \ref{fig:analysis_workflow} depicts our method for analyzing international personal data transfers in Android apps. It departs from the app's package name (a unique identifier for each app in Android, for example, com.companyname.appname) and outputs the app's potential compliance or non-compliance with the GDPR requirements for international transfers. The method consists of four phases: \textit{Application Download}, \textit{Dynamic Analysis}, \textit{Privacy Policy Analysis}, and \textit{Compliance Analysis}. The second and third phases run partially parallel to improve the method's efficiency.

The \textit{Application Download} phase is responsible for obtaining the application's installation file (or APK, which stands for Android Package Kit) and its privacy policy. We download each APK from the Google Play Store through an API \cite{Marty0678/googleplay-api:API}. The privacy policy is obtained by scraping the app's page at the Google Play Store, as it must include a link to the policy text if the app collects personal data.

\begin{figure}
    \centering
    \includegraphics[width=\textwidth]{statics/transfers_outside_EEA.png}
    \caption{International transfers from EEA countries to non-EEA countries. For simplicity, the source of the communication is always shown as the EEA country capital. Destinations legend: United States in grey, United Kingdom in green, Switzerland in purple, Ukraine in yellow, Serbia in blue and Russia in red.}
    \label{fig:transfers_outside_EEA}
\end{figure}

The second phase, \textit{Dynamic Analysis}, runs the app, capturing and analyzing the network traffic generated to spot personal data transfers to anycast addresses. To this end, we first install and launch the app on real mobile devices for two minutes and then automatically stimulate it with the Android Monkey tool (\url{https://developer.android.com/studio/test/other-testing-tools/monkey}) for three minutes. A man-in-the-middle proxy captures the app's network connections, logging the payload sent and the destination IP address, and an instrumentation component deployed in the mobile device logs the app execution traces. We analyze the connections' payloads to retain those transferring personal data off the device. Our approach shares the same limitations as other dynamic analysis tools, for example, incomplete coverage or difficulty in dealing with some user interactions such as login or specific fields. However, it ensures the observation of actual behaviors at scale and therefore offers a view of the lower bound of the phenomenon's occurrence, which in this case is of an enormous scale.

The personal data types detected in the communications include device and user identifiers and location, that is, device model, kernel version and build number; the users' Google Advertisement id; approximate (nearby WiFi BSSIDs) and precise location (connected WiFi BSSID and MAC). All these data types are considered personal data according to the European Union data protection guidelines \cite{JUSTICEwp260rev.01}.

We analyze the connections' destination IPs to keep only anycast addresses and trace them to geolocate their destination. Finally, we analyze the apps' execution traces to identify the code triggering the connection. Since these execution traces preserve the package names (e.g., com. companyname.library), we can identify the TPLs responsible for a captured communication.

Geolocating the destination of a message sent to an anycast address is not a trivial process. By its nature, the geographical origin of the device establishing the connection and the network status determine the destination server and, therefore, its location. We leveraged a previous method \cite{Pascual2024Hunter:Transfers} to determine the destination country of anycast messages, which empirically demonstrated an accuracy of 97.77 percent and a precision of 90.50 percent. The method uses low-level network measurements (i.e., traceroutes) to follow the network packet route and identify the last hop before getting to the destination. This last hop must be close to the destination as it has a direct connection to it, usually in the same city or country. Since these last hops usually have a routing function, they use unicast IP addresses that can be easily geolocated using standard methods (e.g., online geolocation services such as IPInfo).

The destination server of an anycast request depends on the origin of the request; that is, one request sent from Spain will probably get to a different server than another request sent from Finland. Thus, we should geolocate the destination of anycast communications initiated from multiple geographically distributed points to understand the different endpoints comprehensively. Particularly, as we address compliance with GDPR requirements, we must geolocate the destination of anycast communications from all EEA countries. To this end, we leveraged RIPE \cite{HomeAtlas} probes, which offer measurement points distributed throughout the EEA countries. We sent messages from each probe to all the anycast IPs in our list and geolocated the destination, obtaining a map of the source and destination countries for each anycast personal data transfer we spot. We repeated this process 10 times for each source-destination pair so that the results were consistent and we could identify all the possible destinations, considering that the network fluctuations may vary the destination country between executions. The selection of source probes was made by dividing the territory of the EEA countries into cells of 3 degrees longitude and latitude. For each anycast IP, a random probe was selected from each cell, with the condition that it belonged to an EEA country. Approximately 110 RIPE probes were used in each scan, evenly distributed over the whole EEA territory. The use of the RIPE Atlas in this study is only due to the free and convenient geographical distribution of its probes. Using any network measurement platform with similar characteristics would deliver the same results, such as CAIDA ark, PlanetLab, perfSONAR, and so on.

In the third phase, \textit{Privacy Policy Analysis}, automatic analysis of the apps' privacy policy is carried out using machine learning and natural language processing techniques to extract all the information regarding international transfers. Among other data, we obtained the countries disclosed in the privacy policy as possible destinations of personal data transfers. The same process is applied to the TPLs' privacy policies, with the only difference being that one of the researchers manually searched the TPLs' policy text because of the lack of uniformity of their publication.

The fourth and final phase, \textit{Compliance Analysis}, cross-checks the data obtained through our app dynamic analysis and the policy analysis to obtain a conclusion on potential non-compliance issues. On the one hand, the target countries identified for each personal data transfer to an anycast address are triggered by an application and, often, by a TPL it integrates. On the other hand, the set of countries mentioned by the app's (and TPLs if involved) privacy policy as possible destinations for international personal data transfers. After crossing both sets, we obtain applications or TPLs failing to meet GDPR requirements for international transfers of personal data.

Our method is fully automated for retrieving the TPLs' privacy policy. Thus, we have been able to conduct an experiment at a suitable scale to understand the impact of the use of anycast IP addresses by Android apps and, more importantly, by the TPLs they integrate. Although the study has been carried out using probes within the EEA, the method is fully valid in other territories where international personal data transfers are to be studied.

\section{Anycast-Induced Compliance Issues}

\begin{figure}
    \centering
    \includegraphics[width=\textwidth]{statics/transfers_from_EEA_to_UK_and_SW.png}
    \caption{International transfers from EEA countries to anycast servers in United Kingdom (green) and Switzerland (purple).}
    \label{fig:transfers_from_EEA_to_UK_and_SW}
\end{figure}

\begin{figure}
    \centering
    \includegraphics[width=\textwidth]{statics/transfers_from_EEA_to_RU_and_more.png}
    \caption{International transfers from EEA countries to anycast servers in Ukraine (yellow), Serbia (blue) and Russia (red).}
    \label{fig:transfers_from_EEA_to_RU_and_more}
\end{figure}

We randomly selected 6,000 popular (100,000+ downloads) apps from the Google Play Store to analyze their anycast communications and successfully processed 5,759 of them (some experienced issues during the dynamic analysis phase, for example, refused to run on rooted devices). We intercepted 4,278,823 communications from these apps during the experiment, of which 242,497 contained personal data. We further filtered the connections by imposing the conditions that the destination IP address is anycast and the connection contains personal data, thus obtaining the relevant subset of data for the experiment. After this filtering, we have a set of 195,786 intercepted connections, with 200 target anycast IPs addressed from 960 different applications. In addition, our dynamic analysis identified 20 TPLs as triggering 61,443 of these connections.

Having identified the destinations of the international transfers to anycast IP addresses, we retained those transferring personal data outside the EEA. We found that 195 of the 200 IP addresses are involved in personal data transfers outside the EEA, as can be seen in Figs. \ref{fig:transfers_outside_EEA} \ref{fig:transfers_from_EEA_to_UK_and_SW} \ref{fig:transfers_from_EEA_to_RU_and_more}, to the following countries: the United Kingdom (GB), Switzerland (CH), the United States (US), Ukraine (UA), Serbia (RS) and Russia (RU). We cross-checked this data with the applications' privacy policy statements on international transfers and found that 79.58 percent make no mention of international transfers. Furthermore, 98.65 percent of the applications are potentially non-compliant with GDPR in relation to international transfers of personal data as they fail to identify the destination country. The case of TPLs transferring personal data outside the EEA is similar to the apps, as all of them carry out international transfers to anycast IPs, and 90 percent of them potentially breach GDPR by failing to disclose the transfer and the destination country.

All the data generated during our experiments, including a statistical report about the app corpus used, the raw measurements, and the analysis carried out, are available [14], as is the code used to analyze them, at \url{https://github.com/STRAST-UPM/hunter_experiment_anycast_europe}.

\section{Stakeholders' Responsibilities}

Transferring personal data to anycast IP addresses poses a high risk to the citizens' data protection and to the apps' compliance with privacy legislation. Our results show that anycast IP addresses are extensively used to transfer personal data, yet 98.65 percent of apps and 90 percent of TPLs involved are potentially non-compliant with the requirements set by the GDPR.

\begin{figure}
    \centering
    \includegraphics[width=\textwidth]{statics/table_TPLs.png}
    \caption{TPLS and compliance analysis.}
    \label{table:table_TPLs}
\end{figure}

This situation is exacerbated by the integration of TPLs that may automatically collect personal data and send it off the device. The application provider is legally responsible for the processing of personal data in any case, but TPLs should have an adequate privacy policy so that developers are aware of the implications of the use of third-party software. Although the data collection functionalities of TPLs are, in most cases, explicitly triggered by the app developer, there is a significant proportion of cases (9.62 percent) in which the TPL autonomously triggers the connections. We also detected a high aggregation of connections in some TPLs, as can be seen in Table \ref{table:table_TPLs}, indicating that there are companies that have a much higher impact than others. An example of this is Unity Ads (1), which is the source of 57.08 percent of the personal data transfers to anycast addresses we have observed, and the top 5 TPLs by number of connections aggregate 95.88 percent of the connections we have been able to identify. This intensive use of certain TPLs means that any change in their privacy practices will impact a large set of applications and data subjects.

However, obtaining the privacy policies of TPLs is neither straightforward nor a process that can be automated. For the study of the cases detected in the experiment, a systematic but manual search was carried out to locate all privacy policies. The search process consisted of three main steps. The first one involved searching the TPL in the Google Play SDK Index (\url{https://play. google.com/sdks}), where a section called “Guide to the Data Security” sometimes links to the TPL provider. Usually, this link does not lead to a privacy policy but to a website explaining the types of data that can be collected when using the software. Sometimes, the website to which the link in the data security section redirects is just a guide to implementing the software or simply an advertisement. The second step was to locate the TPL in the Maven repository, a widely used tool for managing Android application projects and packages. We could not find any mention of personal data or privacy practices in this software repository, only technical information about versions and installation. Finally, a search was carried out by manually browsing the pages to which the data security section of the Google Play SDK Index redirects or by a direct search for the package name in a web search engine. Using this manual, time-intensive process, we found 80,00 percent of the TPL's privacy policies we searched but could not find the others. This is a concerning result as privacy-related information should be easily accessible to developers in the same way as the technical and implementation information.

The TPLs identified as potentially non-compliant with the legislation are varied, although they belong to three main categories, namely Advertising, Analytics, and Marketing. Advertising TPLs allow developers to monetize their apps by displaying ads to their users. Analytics TPLs collect application usage data as well as possible bugs.

In general, our results show that TPLs are a source of compliance risks regarding international transfers of personal data. Still, there are cases, such as Unity Ads (1) or the Airship SDK (16), where our method did not detect issues. This result is not a coincidence as the above two TPLs are the ones that disclose more countries in their privacy policies as potential destinations of personal data. This contrasts with a common bad practice we found in some privacy policies where it is mentioned that international transfers can be made to any country in the world without mentioning any specific country or mentioning a single country, usually the United States. Marketing TPLs collect data about the users' interaction with the app and its contents, supporting tailoring and delivering marketing campaigns. Many TPLs do not belong to a single category but include features from all three.

We have also detected some interesting practices that could be generalized to support developers in meeting data protection requirements. The case of Appodeal SDK (5) for Android is very interesting because it offers a tool for the generation of privacy policies for any application that includes their library. This kind of tool may help developers to address their data protection obligations. Another practice favoring the proper handling of personal data is one seen in the Official Mixpanel Android SDK (14), which allows constraining personal data transfers only to servers in EEA countries. The decision to use these tools is up to the developer, but it is a best-practice example of how different stake-holders can take responsibility for improving the protection of personal data.

Given the above findings and the critical risks of anycast IP addresses to data protection compliance, this situation must be addressed and corrected. The responsibility of the application providers is evident, as they are responsible for the functionalities they implement and the TPLs they integrate. Nonetheless, they risk severe fines, reputation damage, and degradation of trust. Yet, improvements and controls over TPLs are also necessary to support them in this endeavor, as the ones available do not seem to work properly. Interestingly, the marketplace providers (e.g., Google Play Store) also share their part of the responsibility, for example, in adopting proactive measures for risk assessment and mitigation in order to avoid harm to their users.

Having TPLs declare their privacy practices in a standardized privacy manifest might be a first step in the right direction. This manifest should include at least information on the types of personal data collected, their use, and possible international transfers. As a matter of fact, disclosing the type of data collected and its use is something the industry has already started to implement. For example, both the Apple Store and the Google Play Store provide potential users with a summary of the app's privacy practices, such as through privacy labels. Apple has also implemented a TPL-like manifest in the Apple Store \cite{Third-partyDeveloper}. This manifest includes machine-readable information about the TPL's usage of data. However, none of these initiatives include information on international transfers, which, as we have shown, is a serious concern and a potential data protection risk.

\section{Conclusion}

The use of anycast IP addresses for sending personal data in Android apps is widespread and poses a high risk to personal data protection. This situation is illustrated by the experiment we conducted, where nearly 99 percent of the analyzed apps that send personal data to anycast addresses and 90 percent of the third-party libraries used by these apps are potentially in breach of GDPR requirements for international personal data transfers. While developers are responsible for their apps' compliance with applicable laws, TPLs that collect and transfer personal data should disclose their personal data practices transparently. Unfortunately, the situation we have found is quite the opposite, with no easy access to the TPLs' privacy policies and even not being able to find 20 percent of the privacy policies we thoroughly sought.

Our results call for mechanisms to prevent unintended international transfers and adequately disclose the foreseen ones. To achieve this goal, it is essential that developers have a clear understanding of the TPL's data practices and that they declare them adequately in their apps' privacy policies. In addition to the legal-oriented privacy policy texts, we propose this information to be made available in a standard format that includes all information on the processing of personal data of each TPL, including, of course, information on international transfers. As a matter of fact, this is already starting in the iOS ecosystem \cite{Third-partyDeveloper}, although it does not include any section on international transfers. Although this is a first step in the right direction, it should be extended to other data practices (i.e., international transfers) and ecosystems (i.e., Android), which have been shown to be a serious risk to the privacy and protection of citizens' personal data. At the network level, we are not aware of any technical solution that allows keeping anycast transfers within certain international borders. Our future work is aimed at developing such solutions that can help alleviate the detected situation.

\section{ACKNOWLEDGMENT}

The work of Hugo Pascual was partially supported by the ACES project, funded by the Horizon Europe research program (2021–2027) under grant agreement no. 101093126. The work of Jose M. Del Alamo and David Rodriguez was partially supported by the CEDAR project, funded by the Horizon Europe research program (2021–2027) under grant agreement no. 101135577. The work of Juan C. Dueñas was partially supported by the PRESECREL project, funded by the Plan Estatal de Investigación Científica y Técnica y de Innovación 2017–2020 (Ministerio de Ciencia e Investigación (Spain) - MCIN/AEI/10.13039/501100011033) under Grant agreement PID2021124502OB-C43.

% \bibliographystyle{ieee}
% \bibliography{references}
\printbibliography

\end{document}
